% chirality-time.tex — Section 4: Chirality Requires Time-Orientation
% ASSERT_CONVENTION: natural_units=kB_equals_1, entropy_base=nats, clifford_convention=euclidean_positive_for_Cl6_lorentzian_mostly_minus_for_Cl(d-1,1), complex_structure=u_equals_e7, spin_representation=S10plus_boyle, metric=mostly_minus

%----------------------------------------------------------------------
\section{Chirality Requires Time-Orientation}\label{sec:chirality}
%----------------------------------------------------------------------

Paper~7 showed that a single choice of unit imaginary octonion
$u\in S^6$ determines both the Standard Model gauge group and its
chiral representation~\cite{Paper7}.  That derivation is purely
algebraic: the Euclidean Clifford algebra $\Cl{6}$ decomposes the
$\Spin{10}$ Weyl spinor into left- and right-handed sectors without
reference to spacetime geometry.  We now show that \emph{physical
realization} of this chirality---as spacetime Weyl spinors coupling
to a chiral gauge theory---requires the emergent spacetime to carry a
time-orientation.  This is the new contribution of the present paper:
the algebraic choice $u$ not only builds particle physics structures
(constructive) but also constrains spacetime geometry (constraint).

%----------------------------------------------------------------------
\subsection{The chirality operator in Lorentzian signature}
\label{sec:chirality-lorentz}
%----------------------------------------------------------------------

Consider an even-dimensional Lorentzian manifold of signature
$(d{-}1,1)$, with Clifford generators satisfying
\begin{equation}\label{eq:clifford-lorentz}
  \{\Gamma_\mu,\, \Gamma_\nu\} = 2\,\eta_{\mu\nu}\,,
\end{equation}
where $\eta = \mathrm{diag}(+1,-1,\ldots,-1)$ is the mostly-minus
metric.  In particular, $\Gamma_0^2 = +1$ (timelike) and
$\Gamma_i^2 = -1$ for $i = 1,\ldots,d{-}1$ (spacelike).

The chirality operator (volume form) in $\Cl{d-1,1}$ is
\begin{equation}\label{eq:chirality-op}
  \Gammastar = i^k\,\Gamma_0\,\Gamma_1 \cdots \Gamma_{d-1}\,,
\end{equation}
where the integer $k$ is chosen so that $(\Gammastar)^2 = +1$.
For $d = 4$ this gives the standard $\gamma_5 = i\,\Gamma_0
\Gamma_1\Gamma_2\Gamma_3$ of four-dimensional physics.

The Weyl (chirality) projectors are
\begin{equation}\label{eq:weyl-proj}
  P_L = \frac{1 + \Gammastar}{2}\,,
  \qquad
  P_R = \frac{1 - \Gammastar}{2}\,.
\end{equation}
These are orthogonal idempotents ($P_L + P_R = 1$, $P_L P_R = 0$),
and they decompose the Dirac spinor into Weyl spinors:
$S = S_L \oplus S_R$ with $S_L = P_L S$ and $S_R = P_R S$.

%----------------------------------------------------------------------
\subsection{Time-reversal flips chirality}
\label{sec:time-reversal}
%----------------------------------------------------------------------

Time reversal acts on the Clifford generators as
\begin{equation}\label{eq:time-reversal}
  T\colon\; \Gamma_0 \to -\Gamma_0\,,
  \qquad
  \Gamma_i \to \Gamma_i \quad (i = 1,\ldots,d{-}1)\,.
\end{equation}

\begin{proposition}[Chirality flip]\label{prop:chirality-flip}
For any even $d = 2m$, the chirality operator satisfies
\begin{equation}\label{eq:chirality-flip}
  T(\Gammastar) = -\Gammastar\,.
\end{equation}
\end{proposition}

\begin{proof}
The prefactor $i^k$ is unaffected by~$T$.  The product
$\Gamma_0\Gamma_1\cdots\Gamma_{d-1}$ contains $\Gamma_0$ exactly
once, so under~$T$ it acquires precisely one factor of~$(-1)$:
\begin{equation*}
  \Gamma_0\Gamma_1\cdots\Gamma_{d-1}
  \;\xrightarrow{\;T\;}
  (-\Gamma_0)\Gamma_1\cdots\Gamma_{d-1}
  = -\Gamma_0\Gamma_1\cdots\Gamma_{d-1}\,.
\end{equation*}
Hence $T(\Gammastar) = i^k\cdot(-\Gamma_0\Gamma_1
\cdots\Gamma_{d-1}) = -\Gammastar$.
\end{proof}

\begin{corollary}\label{cor:projector-exchange}
Under time reversal, the Weyl projectors exchange:
\begin{equation}\label{eq:projector-exchange}
  P_L \;\xrightarrow{\;T\;}\; P_R\,,
  \qquad
  P_R \;\xrightarrow{\;T\;}\; P_L\,.
\end{equation}
\end{corollary}

\noindent
The proof is immediate from Eq.~\eqref{eq:chirality-flip}: replacing
$\Gammastar\to -\Gammastar$ in~\eqref{eq:weyl-proj} interchanges
$P_L$ and~$P_R$.

\medskip
We verify the chirality flip explicitly in three dimensions
relevant to this program:

\medskip\noindent
\textbf{$d = 4$} (physical spacetime).
$\gamma_5 = i\,\Gamma_0\Gamma_1\Gamma_2\Gamma_3$,
$(\gamma_5)^2 = +1$\,\checkmark.
Under $T$: $\gamma_5 \to i(-\Gamma_0)\Gamma_1\Gamma_2\Gamma_3
= -\gamma_5$\,\checkmark.

\medskip\noindent
\textbf{$d = 2$} (simplest non-trivial case).
$\Gammastar = \Gamma_0\Gamma_1$,
$(\Gammastar)^2 = +1$\,\checkmark.
Under $T$: $\Gammastar \to -\Gamma_0\Gamma_1
= -\Gammastar$\,\checkmark.

\medskip\noindent
\textbf{$d = 10$} (the $\Spin{10}$ setting of the $\hthree$ program).
$\Gammastar = \Gamma_0\Gamma_1\cdots\Gamma_9$,
$(\Gammastar)^2 = +1$\,\checkmark.
Under $T$: $\Gammastar \to -\Gammastar$\,\checkmark.

\medskip
The result holds for all even-dimensional \emph{Lorentzian}
Clifford algebras $\Cl{d-1,1}$ because $\Gamma_0$ (the unique
timelike generator) appears exactly once in the ordered
product~\eqref{eq:chirality-op}, regardless of~$d$.  In Euclidean
signature there is no timelike direction and the proposition does
not apply.

%----------------------------------------------------------------------
\subsection{Internal vs.\ spacetime chirality}
\label{sec:internal-spacetime}
%----------------------------------------------------------------------

Paper~7's Clifford algebra $\Cl{6}$ and the spacetime Clifford algebra
$\Cl{3,1}$ are distinct objects with distinct chirality operators:

\medskip\noindent
\textbf{Internal chirality} ($\omegasix$, Euclidean).
The volume form of the Euclidean algebra $\Cl{6}$, with generators
$\gamma_1,\ldots,\gamma_6$ satisfying $\{\gamma_i,\gamma_j\} =
2\delta_{ij}$, is~\cite{Paper7}
\begin{equation}\label{eq:omega6-recap}
  \omegasix = \gamma_1\gamma_2\gamma_3\gamma_4\gamma_5\gamma_6\,,
  \qquad
  \omegasix^2 = -1\,,
  \qquad
  (i\,\omegasix)^2 = +1\,.
\end{equation}
The chirality operator $i\,\omegasix$ decomposes the
16-dimensional Weyl spinor $S_{10}^+$ into two 8-dimensional
sectors carrying the Pati-Salam representations
$(\rep{4},\rep{2},\rep{1})$ and
$(\bar{\rep{4}},\rep{1},\rep{2})$.
This decomposition is a \emph{representation-theoretic} fact: it
determines which SM quantum numbers each state carries.  It is
purely algebraic and involves no spacetime geometry.

\medskip\noindent
\textbf{Spacetime chirality} ($\gamma_5$ or $\Gammastar$, Lorentzian).
The chirality operator of $\Cl{3,1}$ (or $\Cl{9,1}$) determines
which Lorentz representation each fermion carries: left-handed
spinors transform as $(\tfrac{1}{2},0)$ under
$\mathrm{SL}(2,\mathbb{C})$, right-handed as $(0,\tfrac{1}{2})$.
This operator involves $\Gamma_0$ and therefore \emph{requires
time-orientation} for a globally consistent definition
(Proposition~\ref{prop:chirality-flip}).

\medskip\noindent
\textbf{The correlation that defines the chiral SM.}
The Standard Model is a chiral gauge theory: its defining feature is
that left-handed fermions couple differently from right-handed ones
under $\SU{2}_L$.  This means a fermion labeled ``left-handed''
must be \emph{simultaneously}:
\begin{enumerate}
  \item in the $(\rep{4},\rep{2},\rep{1})$ representation under
    Pati-Salam (from the internal $\omegasix$), \emph{and}
  \item a left-handed Weyl spinor under the Lorentz group
    (from the spacetime $\gamma_5$).
\end{enumerate}
The internal chirality alone assigns representation labels;
the spacetime chirality determines propagation and coupling.
The full chiral Standard Model requires both chiralities to be
defined and correlated---and the spacetime part requires
time-orientation.

%----------------------------------------------------------------------
\subsection{Conditions for Weyl spinor bundles}
\label{sec:weyl-conditions}
%----------------------------------------------------------------------

We now state precisely which geometric structures a manifold needs
to support Weyl spinor bundles.  A globally consistent
$\Gammastar$ with definite sign requires two independent data:

\begin{enumerate}
  \item \textbf{Space-orientation}: to fix the sign from the ordered
    spatial product $\Gamma_1\cdots\Gamma_{d-1}$.  A permutation of
    two generators flips the sign.

  \item \textbf{Time-orientation}: to fix the sign from $\Gamma_0$.
    Reversing which direction is ``future-pointing'' sends
    $\Gamma_0\to -\Gamma_0$ and hence $\Gammastar\to -\Gammastar$.
\end{enumerate}

\noindent
A Lorentzian spin manifold $(M,g)$ admits Weyl spinor bundles
$S(M) = S_L(M)\oplus S_R(M)$ if and only if $M$ is both
time-oriented and space-oriented~\cite{LawsonMichelsohn1989}.
Conversely, if both orientations are fixed, the ordered frame
$\{e_0,e_1,\ldots,e_{d-1}\}$ (with $e_0$ future-pointing and
$\{e_1,\ldots,e_{d-1}\}$ positively oriented) is globally defined
up to proper orthochronous Lorentz transformations
$\mathrm{SO}^+(d{-}1,1)$, which preserve the sign of~$\Gammastar$.

We assume that the emergent spacetime of Paper~6 provides these
structures:

\medskip\noindent
\textbf{Assumption A6} (continuum limit).
\emph{The continuum limit of Paper~6's self-modeling lattice yields a
smooth Lorentzian manifold $(M,g)$.}

\medskip\noindent
\textbf{Assumption A7} (Lorentzian signature).
\emph{The emergent spacetime has Lorentzian signature, as implied by
Paper~6's causal structure from Lieb-Robinson
bounds~\cite{Paper6}.}

\medskip
Under these assumptions, the lattice structure provides all necessary
geometric data.  A $d$-dimensional regular lattice has $d$ globally
defined basis vectors, which in the continuum limit become a global
frame (a trivialization of the tangent bundle).  A global frame is
\emph{stronger} than each of the conditions it implies:
\begin{equation}\label{eq:framing-hierarchy}
  \text{framing}
  \;\Longrightarrow\;
  \text{spin structure}
  \;\Longrightarrow\;
  \text{orientability}\,,
\end{equation}
where neither implication reverses.%
\footnote{Counterexamples showing strictness: $S^4$ is spin but not
frameable ($\chi(S^4)=2\ne 0$); $\mathbb{C}P^2$ is oriented but
not spin ($w_2\ne 0$).}
Paper~6's lattice additionally provides time-orientation from
Hamiltonian evolution (the sign of $e^{-iHt}$ distinguishes
$t > 0$ from $t < 0$~\cite{Paper6}) and space-orientation from
the lattice ordering.  Therefore all four conditions for Weyl
spinor bundles---spin structure, space-orientation,
time-orientation, and even-dimensionality---are satisfied.

%----------------------------------------------------------------------
\subsection{Three consequences of a single choice}
\label{sec:three-consequences}
%----------------------------------------------------------------------

We are now in a position to extend Paper~7's
``one choice, two consequences'' theorem.  Paper~7 showed that the
single algebraic input $u\in S^6$ determines both the gauge group
and the chiral representation~\cite{Paper7}.
The results of this section add a third consequence: the chirality
from $u$ imposes a geometric constraint on the emergent spacetime.

\begin{theorem}[Three consequences of $u$]
\label{thm:three-consequence}
Let $\hthree$ be the exceptional Jordan algebra with observer
idempotent $E_{11}$.  Assume:
\begin{itemize}
  \item[\textup{(P1)}] $\hthree$ is the universe algebra
    \textup{(Gap~A)}.
  \item[\textup{(P2)}] The observer selects $E_{11}$, breaking
    $\Ffour\to\Spin{9}$ \textup{(Gap~B1)}.
  \item[\textup{(P3)}] A complex structure $u\in S^6 \subset
    \mathrm{Im}(\mathbb{O})$ is chosen \textup{(Gap~B2)}.
  \item[\textup{(P4)}] The complexification principle holds:
    $\Vhalf\otimes_{\mathbb{R}}\mathbb{C} = S_{10}^+$
    \textup{(Gap~C)}.
  \item[\textup{(A6)}] The continuum limit of Paper~6's lattice
    yields a smooth manifold.
  \item[\textup{(A7)}] The emergent spacetime has Lorentzian
    signature.
\end{itemize}
Then the single choice of $u$ simultaneously determines:

\medskip\noindent
\textup{(a)~\textbf{Gauge group}} \textup{(constructive):}\enspace
The $\Ffour$ intersection gives $\SMgauge$,
$\dim = 12$~\cite{Paper7,TodorovDrenska2018}.

\medskip\noindent
\textup{(b)~\textbf{Chirality}} \textup{(constructive):}\enspace
$u$ defines $\Cl{6}\subset\Cl{10}$, whose volume form $\omegasix$
selects the LEFT chiral embedding~\cite{Paper7,Furey2018}:
\begin{equation}\label{eq:chiral-decomp}
  \rep{16}
  \;\to\;
  (\rep{4},\rep{2},\rep{1})
  \;\oplus\;
  (\bar{\rep{4}},\rep{1},\rep{2})\,.
\end{equation}

\medskip\noindent
\textup{(c)~\textbf{Time-orientation requirement}}
\textup{(constraint):}\enspace
The chirality from~\textup{(b)} can only be physically realized as
spacetime Weyl spinors on a time-oriented Lorentzian manifold.
Therefore $u$ creates a necessity for time-orientation on the
emergent spacetime.
\end{theorem}

\begin{proof}
Consequences~(a) and~(b) are proved in Paper~7
(Theorems~4.1 and~4.2 of~\cite{Paper7}).  For~(c), the logical
chain is:
\begin{equation}\label{eq:logical-chain}
  u \;\xrightarrow{\;\text{alg}\;}\; W = u^\perp
  \;\xrightarrow{\;\text{alg}\;}\; \Cl{6}
  \;\xrightarrow{\;\text{alg}\;}\; \omegasix
  \;\xrightarrow{\;\text{rep}\;}\;
  \text{L/R decomp}
  \;\xrightarrow{\;\text{phys}\;}\;
  \text{Weyl spinors}
  \;\xrightarrow{\;\text{geom}\;}\;
  \text{time-oriented}\,.
\end{equation}
Steps~1--4 are pure algebra (from $u$ to the L/R decomposition via
$\omegasix$).  Step~5 is the physical identification of the
algebraic L/R labels with spacetime Weyl spinor types.  Step~6 is
Proposition~\ref{prop:chirality-flip}: the spacetime chirality
operator $\Gammastar$ involves $\Gamma_0$, so a globally consistent
Weyl decomposition requires time-orientation.
The argument is not circular: steps~1--4 assume no geometry, and
time-orientation appears only as a \emph{conclusion}, not a
premise.
\end{proof}

\begin{remark}[Nature of consequence~(c)]\label{rem:constraint-not-constructive}
Consequences~(a) and~(b) are \emph{constructive}: $u$ builds the
gauge group and the chiral representation through explicit algebraic
procedures.  Consequence~(c) is a \emph{constraint}: $u$ imposes a
geometric requirement on the emergent spacetime but does not itself
select a specific time-orientation.  The particular ``arrow of
time''---which direction is the future---comes not from~$u$ but
from the thermodynamic gradient developed in
Section~\ref{sec:gradient-gapc}.
\end{remark}

\begin{remark}[Assumptions A6--A7 are new]\label{rem:new-assumptions}
Assumptions P1--P4 are inherited from Paper~7; they suffice for
consequences~(a) and~(b), which are purely algebraic.
Assumptions~A6--A7 are needed only for consequence~(c) and its
satisfaction: they place the algebraic chirality into a geometric
context.  If A6 or A7 fail, consequences~(a) and~(b) are
unaffected, and the three-consequence theorem reduces to Paper~7's
two-consequence theorem.
\end{remark}

\medskip
That Paper~6's emergent spacetime satisfies the constraint is the
consistency check: under assumptions~A6--A7, the lattice framing
provides a spin structure, the lattice ordering provides
space-orientation, and Hamiltonian evolution provides
time-orientation~\cite{Paper6}.  The chirality from~(b) can
therefore be physically realized.

The key structural insight is that the same choice $u$ that
determines particle physics content (gauge group and chirality) also
constrains spacetime geometry (time-orientation).
This cross-paper link---particle physics from Paper~7 constraining
spacetime from Paper~6---was implicit in the program but has not
been stated before.  It will become central in the next section,
where we show that the time-orientation required by chirality is
the \emph{same} geometric prerequisite needed for the
thermodynamic arrow of time.
